\section*{第 \ref{ch:g1:magnets} 章——磁铁}

\begin{solution}{exo:g1:magnets:1}
被吸住的有：铁钉、回形针、钢勺。未被吸住的有：软木塞、
塑料鸭子、纸片。
\end{solution}

\begin{solution}{exo:g1:magnets:2}
铁（以及主要成分为铁的钢）。无论如何都不被\omterm{def:g1:magnets:magnet}{吸引}的有：硬币和
厨房锡箔纸 —— 都是金属，但不是铁。
\end{solution}

\begin{solution}{exo:g1:magnets:3}
纸片（画）夹在它们中间，\omterm{def:g1:magnets:magnet}{磁铁}的
拉力直接穿过纸片吸到铁门上，它能
穿过物体起作用。
\end{solution}

\begin{solution}{exo:g1:magnets:4}
把\omterm{def:g1:magnets:magnet}{磁铁}紧贴玻璃杯的\emph{外侧}，靠近
回形针，慢慢沿杯壁向上移动：回形针在杯内跟着移动，
直到杯口，被磁力穿过玻璃和水拉着上升。
\end{solution}

\begin{solution}{exo:g1:magnets:5}
离\omterm{def:g1:magnets:magnet}{磁铁}较远时，拉力太弱，无法移动回形针；只有
靠得很近时拉力才变强。\omterm{def:g1:magnets:magnet}{磁铁}的作用范围随距离增大而迅速减弱。
\end{solution}

\begin{solution}{exo:g1:magnets:6}
把两块\omterm{def:g1:magnets:magnet}{磁铁}中的一块调转方向，使另一端
先靠近。这时两块\omterm{def:g1:magnets:magnet}{磁铁}会互相排斥，而不是吸到一起。
\end{solution}

\begin{solution}{exo:g1:magnets:7}
在两端——回形针基本不受中间部分的\omterm{def:g1:magnets:magnet}{吸引}。\omterm{def:g1:magnets:magnet}{磁铁}在
两端磁性最强。
\end{solution}

\begin{solution}{exo:g1:magnets:8}
相同点：两者都能把物体拉向你。不同点：玩具必须用手
碰到才能拉动，而\omterm{def:g1:magnets:magnet}{磁铁}可以隔着空气、
木头和水拉动回形针，无需接触。
\end{solution}

\begin{solution}{exo:g1:magnets:9}
依次把\omterm{def:g1:magnets:magnet}{磁铁}靠近每个物体。铁钉会被吸过去；硬币
和锡箔纸虽然都是金属，却不动。所以汤姆的
说法不对；改正为：“\omterm{def:g1:magnets:magnet}{磁铁}能\omterm{def:g1:magnets:magnet}{吸引}由铁（或钢）制成的物体，
而不是所有金属。”
\end{solution}
